\documentclass{article}

\usepackage{fancyhdr}
\usepackage[letterpaper,margin=0.5in,includehead,includefoot,headheight=14pt,headsep=6pt]{geometry}
\usepackage{microtype}
\usepackage{multicol}
\usepackage{enumitem}
\usepackage{amsmath,amssymb}
\usepackage[hidelinks]{hyperref}

\setlength{\columnsep}{0.22in}
\setlength{\parindent}{0pt}
\setlength{\parskip}{3pt}
\setlist{nosep,leftmargin=*}
\setlength{\abovedisplayskip}{3pt}
\setlength{\belowdisplayskip}{3pt}

\pagestyle{fancy}
\fancyhf{}
\fancyhead[L]{\textbf{Counting Core}}
\fancyhead[R]{Name: \rule{2.2in}{0.4pt}}
\fancyfoot[C]{\thepage}
\renewcommand{\headrulewidth}{0.4pt}

\begin{document}
\small

\textbf{Public practice bank.} These pages are not a live graded form and are
not themselves scored. A supervised attempt lasts 20 minutes and contains
exactly five problems, one chosen from each section A--E. Every lettered part
under a problem label belongs to that one problem. \textbf{Meets} requires four
of the five problems to be fully correct; \textbf{Exceeds} requires all five.
Staff records the result in Canvas as 0 (Not Yet), 4 (Meets), or 5 (Exceeds).
Every retake is a new version of the entire five-problem Core.

\textbf{Directions.} Show enough work to identify and justify the counting or
probability model you used. Exact fractions and unsimplified products are
acceptable unless a decimal is requested.

\section*{A. Fundamental Counting Rules}

% Problem families adapted from Rosen 7e, Sections 6.1--6.2.
\begin{enumerate}[label=\textbf{A.\arabic*.}]
    \item A cafe offers 5 soups and 4 salads. A lunch includes either one soup
    or one salad, but not both. How many lunches are possible? Name the counting
    rule.

    \item A password contains two uppercase letters followed by three digits.
    Repetition is allowed. How many passwords are possible? Give an unsimplified
    product and name the counting rule.

    \item Let $A=\{1,2,3,4,5\}$ and $B=\{1,2,\ldots,8\}$. How many functions
    $f:A\to B$ exist? Identify the sequence of choices being counted.

    \item If $|A|=19$, $|B|=20$, and $|A\cap B|=8$, find $|A\cup B|$. Write
    the formula you used and name the principle.

    \item A hash table maps 40 distinct keys to 20 possible hash values. What
    bucket occupancy is guaranteed? State the generalized pigeonhole calculation
    and explain why the data do not guarantee a larger occupancy.
\end{enumerate}

\section*{B. Arrangements and Selections}

% Problem families adapted from Rosen 7e, Sections 6.3--6.4.
\begin{enumerate}[label=\textbf{B.\arabic*.}]
    \item A committee of 3 students is selected from 10 students. How many
    committees are possible? Give a binomial-coefficient expression and explain
    why order does not matter.

    \item A president and a vice president are selected from 10 students, with
    no student holding both offices. How many outcomes are possible? Give a
    permutation or product expression and explain why order matters.

    \item How many distinct strings can be made by reordering the letters of
    \texttt{BANANA}? Give the appropriate multinomial expression.

    \item How many bit strings of length 8 contain exactly three $1$s? Give a
    binomial-coefficient expression and identify what is being selected.

    \item Find the coefficient of $x^3$ in $(1+x)^8$. State the binomial-theorem
    term that determines your answer.
\end{enumerate}

\newpage

\section*{C. Constrained Counting}

% Problem families adapted from Rosen 7e, Sections 6.1, 6.3, and 6.4.
\begin{enumerate}[label=\textbf{C.\arabic*.}]
    \item How many 5-card hands from a standard 52-card deck contain exactly
    two hearts? Give an unsimplified product of binomial coefficients.

    \item How many bit strings of length 8 contain at least two $1$s? Use a
    complement and show the subtracted cases.

    \item A group of 5 students is selected from 12 computer-science majors and
    8 mathematics majors. How many groups contain exactly 3 computer-science
    majors? Give an unsimplified product of binomial coefficients.

    \item Let $|A|=5$ and $|B|=8$. How many injective functions $f:A\to B$
    exist? Give a permutation or product expression and explain the restriction.

    \item A six-character code uses uppercase letters without repetition and
    must begin with a vowel. How many codes are possible? Treat
    $A,E,I,O,U$ as the five vowels and give an unsimplified product.
\end{enumerate}

\section*{D. Probability Models}

% Problem families adapted from Rosen 7e, Sections 7.1--7.2.
\begin{enumerate}[label=\textbf{D.\arabic*.}]
    \item A fair coin is tossed and a fair six-sided die is rolled. Find the
    probability that the coin is heads or the die is even (or both). Show the
    favorable-outcome count or an equivalent probability calculation.

    \item Suppose $P(A)=0.55$, $P(B)=0.40$, and $P(A\cap B)=0.25$. Find
    $P(A\cup B)$ and $P(A^c\cap B)$, showing the probability identities used.

    \item A uniformly chosen 8-bit string is selected. Find the probability
    that it contains exactly three $1$s. Give an exact fraction.

    \item A bag contains 4 red and 6 blue tokens. Two tokens are drawn without
    replacement. Find the probability that the first is red and the second is
    blue. Show the conditional product.

    \item In a class, 24 students use Python, 18 use Java, and 10 use both.
    One student is chosen uniformly. Find
    $P(\text{Python}\mid\text{Java})$ and
    $P(\text{Java}\mid\text{Python})$.
\end{enumerate}

\newpage
\thispagestyle{fancy}

\section*{E. Modeling and Interpretation}

% Problem families adapted from Rosen 7e, Sections 7.2--7.4.
\begin{enumerate}[label=\textbf{E.\arabic*.}]
    \item Events $A$ and $B$ satisfy $P(A)=0.6$, $P(B)=0.5$, and
    $P(A\cap B)=0.3$. Determine whether $A$ and $B$ are independent, showing
    the numerical test, and then find $P(A\mid B)$.

    \item A condition occurs in $1\%$ of a population. A test is positive for
    $95\%$ of people who have the condition and for $8\%$ of people who do not.
    Find $P(\text{positive})$ and then
    $P(\text{condition}\mid\text{positive})$. Give a one-sentence base-rate
    interpretation.

    \item Factory A supplies $40\%$ of a store's batteries and $3\%$ of its
    batteries are defective. Factory B supplies the rest and $1\%$ of its
    batteries are defective. A randomly selected battery is defective. Find the
    probability that it came from Factory A. Show total probability and Bayes'
    rule.

    \item A game uses a fair six-sided die. It pays \$10 for a 6, \$2 for a 4
    or 5, and \$0 otherwise. A play costs \$2. Find the expected \emph{net}
    value of one play and interpret the sign.

    \item The table gives successes/participants for two treatments.
    \[
    \begin{array}{c|cc}
      &\text{Treatment A}&\text{Treatment B}\\\hline
      \text{Group 1}&30/100&400/1000\\
      \text{Group 2}&800/1000&90/100
    \end{array}
    \]
    Identify the better treatment within each group and overall. Then explain in
    one sentence how the unequal treatment denominators produce the reversal.
\end{enumerate}


\end{document}
